\documentclass{article}
\usepackage{amsthm}
\newtheorem{theorem}{Theorem}
\newtheorem{definition}{Definition}
\begin{document}
\begin{definition}
A metric space is \emph{homogeneous} if its isometry group acts transitively on points.
\end{definition}
\begin{definition}
A \emph{Kripke frame} is an ordered pair $(W,R)$ where $W$ is nonempty and $R$ is a relation on $W$.
\end{definition}
\begin{definition}
A subgroup $H$ of a group $G$ is called \textbf{malnormal} if $gHg^{-1}\cap H$ is trivial for every $g$ outside $H$.
\end{definition}
\begin{definition}
A group $G$ is called a \emph{Tarski monster} if every proper nontrivial subgroup of it has prime order.
\end{definition}
\begin{theorem}
Let $X$ be a homogeneous metric space. Then $X$ is a metric space.
\end{theorem}
\begin{theorem}
Let $X$ be a compact metric space. If $X$ is homogeneous, then $X$ is complete.
\end{theorem}
\begin{theorem}
Let $F$ be a Kripke frame. Then $F=F$.
\end{theorem}
\begin{theorem}
Let $G$ be a group and let $H$ be a malnormal subgroup of $G$. Then $H$ is a subgroup of $G$.
\end{theorem}
\begin{theorem}
Every Tarski monster is a group. Let $G$ be a Tarski monster. Then $G$ is nontrivial.
\end{theorem}
\begin{definition}
A loop $L$ is said to be \emph{Moufang} if it satisfies the Moufang identities, and a Moufang loop $L$ is called a \emph{code loop} if it has a central subgroup of order two with an elementary abelian quotient.
\end{definition}
\begin{definition}
If the generating set of a cone is linearly independent, then the cone is said to be \emph{simplicial}. Such a group is called a \emph{monster}. Finite-dimensional algebras $A$ with this property are called \emph{weakly Gorenstein} algebras.
\end{definition}
\begin{theorem}
Let $L$ be a Moufang loop and let $K$ be a code loop. Then $L=L$ and $K=K$.
\end{theorem}
\begin{theorem}
Let $G$ be a monster. Then $G$ is a group.
\end{theorem}
\begin{theorem}
Let $C$ be a simplicial cone. Then $C=C$. Let $A$ be a finite-dimensional algebra. If $A$ is weakly Gorenstein, then $A=A$.
\end{theorem}
\end{document}
